\specialsection{\texorpdfstring{\((\mathbf{R}/a\mathbf{Z})^{n}\) 中的公式表}{(R/aZ)\^{}n 中的公式表}}

为了读者方便，我们对群 \((\mathbf{R}/a\mathbf{Z})^{n}\)（其中 a \textgreater{} 0，赋以总质量为 1 的 Haar 测度，记作 \(a^{-n} dx\)）叙述 Fourier 理论的主要公式。当 n = 1 时，这对应于周期为 a 的周期函数。\((\mathbf{R}/a\mathbf{Z})^{n}\) 的对偶群与群 \(Z^{n}\) 等同（通过映射 \((\mathbf{R}/a\mathbf{Z})^{n} \times \mathbf{Z}^{n} \to \mathbf{U}\)），该映射是由下列映射过渡到商而诱导的

\[
(x, y) \mapsto \exp \left(\frac {2 i \pi}{a} x \cdot y\right)
\]

（对所有 \(x \in R^{n}\) 以及每个 \(y \in Z^{n}\)）。\(Z^{n}\) 上的对偶测度于是就是计数测度。

\[
\begin{array}{l l} f \in \mathrm{L} ^ {1} ((\mathbf {R} / a \mathbf {Z}) ^ {n}) & \widehat {f} (y) = \frac {1}{a ^ {n}} \int_ {(\mathbf {R} / a \mathbf {Z}) ^ {n}} f (x) \exp \Bigl (- \frac {2 i \pi}{a} y \cdot x \Bigr) d x \\ & (y \in \mathbf {Z} ^ {n}; \text {Def. 3, p. 206}) \\ f \in \mathrm{L} ^ {1} ((\mathbf {R} / a \mathbf {Z}) ^ {n}) \text {or} & \widehat {f} _ {h} (y) = \exp (\frac {2 i \pi}{a} h \cdot y) \widehat {f} (y) \\ \mathrm{L} ^ {2} ((\mathbf {R} / a \mathbf {Z}) ^ {n}) & f _ {h} (x) = f (x + h) \\ h \in (\mathbf {R} / a \mathbf {Z}) ^ {n} & (\text {Eq. (11), p. 208}) \\ f \in \mathrm{L} ^ {1} ((\mathbf {R} / a \mathbf {Z}) ^ {n}) \text {or} & \widehat {g} _ {h} (y) = \widehat {f} (y - h) \\ \mathrm{L} ^ {2} ((\mathbf {R} / a \mathbf {Z}) ^ {n}) & g _ {h} (x) = f (x) \exp (\frac {2 i \pi}{a} h \cdot x) \\ h \in (\mathbf {R} / a \mathbf {Z}) ^ {n} & (\text {Eq. (12), p. 208}) \\ f \in \mathrm{L} ^ {2} ((\mathbf {R} / a \mathbf {Z}) ^ {n}) & \frac {1}{a ^ {n}} \int_ {(\mathbf {R} / a \mathbf {Z}) ^ {n}} | f (x) | ^ {2} d x = \sum_ {y \in \mathbf {Z} ^ {n}} | \widehat {f} (y) | ^ {2} \\ & (\text {Th. 1, p. 215}) \\ f \in \mathrm{L} ^ {2} ((\mathbf {R} / a \mathbf {Z}) ^ {n}) \text {such that} & f (x) = \sum_ {y \in \mathbf {Z} ^ {n}} \widehat {f} (y) \exp (\frac {2 i \pi}{a} y \cdot x) \\ \widehat {f} \in \mathrm{L} ^ {1} (\mathbf {Z} ^ {n}) & (\text {almost everywhere; Prop. 12, p. 219}) \\ f \in \mathrm{A} ((\mathbf {R} / a \mathbf {Z}) ^ {n}) & f (x) = \sum_ {y \in \mathbf {Z} ^ {n}} \widehat {f} (y) \exp (\frac {2 i \pi}{a} y \cdot x) \\ & (x \in (\mathbf {R} / a \mathbf {Z}) ^ {n}; \text {Prop. 11 p. 217}) \end{array}
\]

